\input{\string~/Documents/School/"UC Irvine"/ta_template2.0.tex}
\rh{2B}{7.8}
\chead{\today}
\graphicspath{ {\string~/Desktop} }
\usetikzlibrary{arrows}
%============ANSWER KEY TOGGLE===========
\newcommand{\ans}[1]{ \noindent {\color{red} \textbf{Answer:} #1}}
\renewcommand{\ans}[1]{\vspace{3in}}
%==========================================
\begin{document}
\begin{center}
\textbf{Improper Integrals}
\end{center}

\begin{aun}[Improper Integrals]
The big idea is that one can approximate infinite-area integrals with limits of finite ones. We'll summarize these into three types:
\begin{description}
\item[Case 1: Infinite Bounds] \[
\int_a^\infty f(x) dx = \lim_{b \to \infty} \int_a^b f(x)dx
\]
where the integral converges if the corresponding limit does, and diverges if the corresponding limit diverges. We can do this with one sided limits for discontinuities also. 
\item[Case 2: Discontinuous Integrands] When $f(x)$ is discontinuous at some $c$ with $a < c < b$, we write:
\[
\int_a^b f(x)dx = \lim_{e \to c^-} \int_a^e f(x) dx + \lim_{g \to c^+} \int_g^b f(x) dx
\]
and say that the integral \emph{converges} if each of these limits converges when studied individually. 
\end{description}
\end{aun}
\begin{prop}
We have that:
\[
\int_{1}^\infty \frac{1}{x^p} dx \begin{cases}
\text{ converges if } & p > 1\\
\text{ diverges if } & p \leq 1
\end{cases}
\]
\end{prop}
\begin{thm}[Comparison Theorem]
Suppose that $f$ and $g$ are continuous functions with $f(x) \geq g(x) \geq 0$ for all $x \geq a$:
\begin{enumerate}[label=\alph*.]
\item If $\int_a^\infty f(x) dx$ is convergent, then $\int_a^\infty g(x) dx$ is convergent.
\item If $\int_a^\infty g(x) dx$ is divergent, then $\int_a^\infty f(x) dx$ is divergent.
\end{enumerate}
\end{thm}


\begin{exx}
Determine whether the following integrals are convergent or divergent:
\[
\int_{3}^\infty \frac{1}{(x-2)^{3/2}} dx \mte{ and } \int_{2}^\infty e^{-5p} dp
\]
\end{exx}
\ans{For the first integral, we can make the substitution $u = x - 2$ and notice $du = dx$; further, 
\[
\int_3^\infty \frac{1}{(x-2)^{3/2}} = \int_1^\infty \frac{1}{u^{3/2}} du
\]
which converges by the $p$ test. For the second integral, we first have to take an antiderivative of the integrand:
\[
\int_{2}^\infty e^{-5p} dp = \lim_{b \to \infty}  \restr{\frac{e^{-5p}}{-5}}{2}^b = \lim_{b \to \infty} \frac{-1}{5} \lrb{\frac{1}{e^{5b}} - \frac{1}{e^{10}}}
\]
which surely converges:
\[
\int_{2}^\infty e^{-5p} dp = \frac{1}{5e^{10}}
\] 
}




\problembox{Determine the convergence of the following integrals:
\begin{enumerate}[label=\alph*.]
\item $\displaystyle \int_0^\infty \frac{x^2}{\sqrt{1+x^3}}dx$
\item $\dst \int_1^\infty \frac{1}{x^2 + x}dx$
\item $\dst \int_0^4 \frac{dx}{x^2 - x - 2}$
\end{enumerate}
}

\ans{
For $(a)$, we can make the substitution: $u = 1 + x^3$, $du = 3x^2 dx$
\[
\int_0^\infty \frac{x^2}{\sqrt{1+x^3}}dx = \lim_{b \to \infty} \frac13 \int_1^\infty \frac{du}{\sqrt{u}}
\]
Without computation, this diverges by the $p$-test. \\
The second integral should suggest comparison; since we know that increasing the denominator decreases the function, we can bound $\frac{1}{x^2 + x} \leq \frac{1}{x^2}$ which we know converges by the $p$-test. Determining the actual value of this integral is quite challenging, but classifying whether or not it does is important. \\
We show $(c)$. We have a polynomial in the denominator, so we factor and use partial fraction decomposition. 
\[
\frac{1}{x^2 - x - 2} = \frac{1}{3(x-2)} - \frac{1}{3(x+1)}
\]
This means that our function has discontinuities at $x=-1$ and $x= 2$. Because only $x =2$ is in the range of our integral, this is a case 2 integral:
\[
\lim_{b \to 2^-} \int_0^b \frac{1}{x^2 - x + 2} dx + \lim_{b \to 2^+} \int_b ^4 \frac{1}{x^2 - x + 2} dx
\]
The remainder of the problem uses the partial fraction decomposition to compute these integrals. This integral \emph{does not converge.} (Why? eliminate the integrals and study the resulting limits. )
}

\problembox{[Some More Twists] Consider these slightly more challenging problems as well!
\begin{enumerate}[label=\alph*.]
\item Determine the values of $p$ for which the integral converges in:
\[
\int_0^1 \frac{1}{x^p} dx 
\]
\item $\dst \int_e^\infty \frac{1}{x (\ln (x))^2} dx$
\item $\dst \int_0^1 \frac{dx}{\sqrt{1-x^2}}$
\end{enumerate}
}

\ans{(a) can be solved by just following the exact steps that we have been doing before- this function has a discontinuity at $0$, so we evaluate:
\[
\lim_{b \to 0^+} \int _0^1 x^{-p} = \lim_{b \to 0^+} \restr{\frac{x^{1-p}}{1-p}}{b}^1 = \lrb{\frac{1}{1-p} - \lim_{b \to 0^+} \frac{b^{1-p}}{1-p}}
\]
converges to $\frac{1}{1-p}$ when $p < 1$. \\
(b) can be done with a direct substitution, this is an integral we have studied before. \\
(c) can be done by recognizing it as a trigonometric substitution, this \emph{does} converge!
}

\end{document}
